\newcommand{\DoNotLoadEpstopdf}{}
\documentclass[11pt]{article}
\usepackage[T1]{fontenc}
\usepackage{lmodern}
\usepackage[margin=1.08in]{geometry}
\usepackage{amsmath,amssymb,amsthm,mathtools}
\usepackage{microtype}
\usepackage{xcolor}
\usepackage[colorlinks=true,linkcolor=blue!45!black,citecolor=blue!45!black,urlcolor=blue!45!black]{hyperref}
\input{glyphtounicode}
\input{glyphtounicode-cmex}
\pdfgentounicode=1
\pdftrailerid{}
\pdfsuppressptexinfo=15
\hypersetup{pdftitle={Graph Decompositions at the Integral Expectation Threshold},pdfauthor={OpenAI}}
\newtheorem{theorem}{Theorem}[section]
\newtheorem{lemma}[theorem]{Lemma}
\newtheorem{proposition}[theorem]{Proposition}
\newtheorem{corollary}[theorem]{Corollary}
\theoremstyle{remark}
\newtheorem{remark}[theorem]{Remark}
\numberwithin{equation}{section}
\newcommand{\E}{\mathbb E}
\newcommand{\Prob}{\mathbb P}
\newcommand{\calD}{\mathcal D}
\newcommand{\calF}{\mathcal F}
\newcommand{\calQ}{\mathcal Q}
\newcommand{\calR}{\mathcal R}
\DeclareMathOperator{\degen}{degen}
\DeclareMathOperator{\Bin}{Bin}
\title{Graph Decompositions at the Integral Expectation Threshold}
\author{OpenAI}
\date{October 5, 2026}

\begin{document}
\maketitle

\begin{abstract}
We prove the graph-decomposition conjecture of Ascoli, He, Park,
and Talagrand. Every graph admits a partition into a universally
bounded number of fixed edge pieces, each having ordinary containment
threshold at most a universal constant times the original graph's
integral expectation threshold. The partition is chosen before
sampling the random host, and the separate embeddings of the pieces
need not agree on shared vertices.
\end{abstract}

\section{Introduction}
\label{sec:intro}

All graphs are finite and simple, and copies are not required to
be induced. For a graph $J$ on at most $n$ vertices, let
\[
 p_c(J;n)=\min\{p\in[0,1]:
           \Prob(G(n,p)\text{ contains a copy of }J)\geq1/2\}.
\]
Here $G(n,p)$ has vertex set $[n]=\{1,\ldots,n\}$, with edges
present independently with probability $p$. We set $p_c(J;n)=0$
when $J$ is edgeless.

To define the integral expectation threshold, put $X=E(K_n)$ and
let $\calF_H\subseteq2^X$ be the increasing family of edge sets
whose graphs contain $H$. A family $\mathcal C\subseteq2^X$
\emph{covers} $\calF_H$ if every $A\in\calF_H$ contains some
$S\in\mathcal C$. Its cost at $p$ is
\[
 c_p(\mathcal C)=\sum_{S\in\mathcal C}p^{|S|}.
\]
For $H$ with at least one edge, define
\begin{equation}\label{eq:q-definition}
 q(H;n)=\max\{p\in[0,1]:
    \calF_H\text{ has a cover }\mathcal C
    \text{ with }c_p(\mathcal C)\leq1/2\}.
\end{equation}
The maximum exists because there are only finitely many families
of subsets of $X$. Each cover gives a first-moment bound on
containment, so $q(H;n)\leq p_c(H;n)$. For convenience set
$q(H;n)=0$ for edgeless graphs as well.

\begin{theorem}\label{thm:main}
There exist an absolute integer $k\geq2$ and an absolute constant
$L\geq1$ such that, for every $n\geq2$ and every graph $H$ on
at most $n$ vertices, the edges of $H$ admit a partition
\[
 E(H)=E(H_1)\,\dot\cup\,\cdots\,\dot\cup\,E(H_k)
 \quad\text{with}\quad
 p_c(H_i;n)\leq Lq(H;n)\qquad(1\leq i\leq k).
\]
The pieces may be empty and may share vertices. The partition
depends only on $H,n$ and is fixed before the random host is sampled.
\end{theorem}

This proves Conjecture~8.1 of Ascoli, He, Park, and Talagrand
\cite{AHPT26}. Their formulation allows overlapping edge sets.
Assigning each edge to one piece containing it gives a partition
and cannot increase any piece's containment threshold. Thus the
two formulations are equivalent.

\subsection*{History and significance}
Kahn and Kalai \cite{KK07} proposed a general comparison between
ordinary thresholds and integral expectation thresholds with a
logarithmic loss. Following the fractional version established by
Frankston, Kahn, Narayanan, and Park \cite{FKNP21}, Park and Pham
\cite{PP24} proved the integral comparison. For graph containment
it gives
\[
 p_c(H;n)\leq Cq(H;n)\log\max\{2,|E(H)|\}
\]
with an absolute constant $C$. Theorem~\ref{thm:main} removes
this logarithmic loss after partitioning the target into a bounded
number of pieces. The assertion concerns each fixed piece's
ordinary threshold; their embeddings need not agree on shared
vertices.

Ascoli, He, Park, and Talagrand introduced generalized thresholds
based on unions of likely graphs, related them to Talagrand's
discrete-convexity conjecture, and formulated the decomposition
problem above \cite{AHPT26}. They proved bounded-degeneracy
cases and further cases with both degeneracy and maximum-degree
restrictions \cite[Theorems~1.12--1.13]{AHPT26}. Their
Section~8 isolates the bipartite edges between high- and
low-degree vertices as a principal remaining difficulty.

We use the same-parameter integral covering theorem from
\emph{Talagrand's discrete-convexity conjecture}
\cite[Theorem~1.1]{OpenAI26DC}, which proves the discrete
conjecture formulated in \cite[Conjecture~7.1]{Talagrand10}.
It is stated precisely in Section~\ref{sec:prelim}; this is an
essential companion input. All remaining graph-decomposition and
embedding arguments are proved below.

\subsection*{Method}
The covering theorem places one labelled copy of $H$ inside the
union of a bounded number of deterministic graphs from a class
$\calD$ having high probability at a density $r$ comparable to
$q(H;n)$. Assigning each edge to one of these graphs gives initial
pieces. In the large-degeneracy case we isolate a small set $Y$
containing the high-degree vertices and enlarge it until each
outside vertex has few neighbors in $Y$. Subsequent splitting
produces pieces inside $Y$, outside $Y$, and across the cut.

To embed a cut piece while fixing $Y$, each nonisolated outside
vertex needs a distinct image adjacent to its prescribed set of
at most $s$ vertices in $Y$. Hall's criterion asks for enough
candidate images for every collection of outside vertices.
Membership in a typical class alone does not ensure this; the
class must control unions of arbitrarily many such conditions.

The compression lemma selects at most $\lceil2r^{-s}\rceil$
requirements from any family so that every neighborhood containing
an original requirement contains a selected requirement with
probability at least $1/2$ after adding each vertex of $Y$
independently with probability $r$. This principle is independent
of graphs. It permits simultaneous concentration over all
subfamilies of this bounded size while still controlling every
Hall condition. We define $\calD$ using these tests for all small
$Y$ at once, so that transfer applies to cut pieces chosen after
the deterministic cover.

For the two other regions, a direct embedding lemma uses successive
matchings of layers of pairwise nonadjacent vertices. Graphs of
bounded degeneracy are handled by small star forests, adapting
the degeneracy splitting and center-testing arguments of
\cite[Claim~3.9 and Section~6]{AHPT26}.

Section~\ref{sec:prelim}
collects the covering input and elementary estimates.
Section~\ref{sec:transfer} proves neighborhood transfer,
Section~\ref{sec:direct} gives the direct embedding lemmas,
and Section~\ref{sec:assembly} constructs the partition and
proves Theorem~\ref{thm:main}.

\section{Covering input and elementary estimates}
\label{sec:prelim}

The first step converts a family of typical labelled graphs into
a deterministic cover of one copy of the target. We then record
the elementary degree and threshold bounds used to subdivide that
copy.

For a finite ground set $X$, write $\mu_r$ for the law of its
Bernoulli-$r$ random subset. Covering and the cost $c_r$ have the
same meaning as in the introduction, with $X$ in place of $E(K_n)$.

\begin{theorem}[Discrete convexity, {\cite[Theorem~1.1]{OpenAI26DC}}]
\label{thm:discrete}
Set $K_0=2^{75}$. For every nonempty finite set $X$, every
$r\in(0,1)$, and every family $\calD\subseteq2^X$ with
$\mu_r(\calD)\geq1-1/K_0$, the family
\[
 \mathcal E_{K_0}(\calD)=
 \{S\subseteq X:S\not\subseteq D_1\cup\cdots\cup D_{K_0}
                   \text{ for every }D_1,\ldots,D_{K_0}\in\calD\}
\]
has a cover of cost at most $1/2$ at the same parameter $r$.
The family $\calD$ need not be monotone, and members of a tuple
may repeat.
\end{theorem}

\begin{corollary}\label{cor:typical-cover}
Let $n\geq2$, let $H$ have an edge and at most $n$ vertices,
and suppose $q(H;n)<r<1$. If a family of graphs on $[n]$
satisfies $\mu_r(\calD)\geq1-1/K_0$, then some copy of $H$
has its edges contained in the union of $K_0$ members of $\calD$.
\end{corollary}
\begin{proof}
If no such copy exists, no graph containing $H$ can be contained
in such a union. Thus $\calF_H\subseteq\mathcal E_{K_0}(\calD)$.
Theorem~\ref{thm:discrete} supplies a cover of $\calF_H$ of
cost at most $1/2$ at $r$, contrary to $r>q(H;n)$.
\end{proof}

For a graph $J$, let
\[
 V_+(J)=\{v\in V(J):\deg_J(v)>0\}
\]
be its set of nonisolated vertices. Isolated vertices can be
deleted during an embedding and added afterward: any embedding
of the remaining vertices extends injectively when $|V(J)|\leq n$.
If $J\subseteq H$, then $\calF_H\subseteq\calF_J$, so every
cover of $\calF_J$ covers $\calF_H$. In particular,
\begin{equation}\label{eq:monotonicity}
 q(J;n)\leq q(H;n),\qquad p_c(J;n)\leq p_c(H;n).
\end{equation}

The \emph{degeneracy} $d=\degen(H)$ is the maximum of the
minimum degrees of nonempty subgraphs of $H$. Equivalently,
the vertices admit an order with at most $d$ preceding neighbors
per vertex. To obtain the order, repeatedly remove a vertex of
degree at most $d$ and place it last. Conversely, the last vertex
of any subgraph in such an order has degree at most $d$.
Counting each edge at its later endpoint shows that every
subgraph $J$ of $H$ satisfies
\begin{equation}\label{eq:edge-bound}
 |E(J)|\leq d|V(J)|.
\end{equation}

\begin{lemma}\label{lem:q-lower}
If $H$ has an edge, $|V(H)|\leq n$, and $d=\degen(H)$, then
\begin{equation}\label{eq:q-lower}
 q(H;n)\geq\frac1{2\binom n2},\qquad
 q(H;n)\geq\frac12 n^{-2/d}.
\end{equation}
\end{lemma}
\begin{proof}
The family of all single edges covers $\calF_H$, proving the
first bound. For the second, choose a subgraph of minimum degree
$d$, with $v$ vertices and $a\geq dv/2$ edges. Its copies on
$[n]$ form a cover of $\calF_H$ with cost at most $n^vp^a$.
At $p=2^{-1/a}n^{-v/a}\in(0,1)$ this is at most $1/2$.
Since $a\geq1$ and $v/a\leq2/d$, this parameter is at least
$\frac12n^{-2/d}$.
\end{proof}

We use the following elementary splitting from
\cite[Claim~3.9]{AHPT26}.

\begin{lemma}[Degeneracy splitting]\label{lem:split}
For every positive integer $m$, the edges of a $d$-degenerate
graph can be partitioned into $m$ graphs, each of degeneracy
at most $\lceil d/m\rceil$.
\end{lemma}
\begin{proof}
Fix a degeneracy order. At each vertex, partition its edges to
preceding neighbors into $m$ groups of size at most
$\lceil d/m\rceil$, and assign each group to its corresponding
piece. Each edge is assigned once. The same order certifies the
claimed bound in every piece.
\end{proof}

\section{Uniform neighborhood transfer}
\label{sec:transfer}

We construct a typical class whose small-neighborhood bipartite
subgraphs embed in a slightly denser random graph. To embed a
bipartite graph $J$ while fixing one part $Y$, each nonisolated
vertex $w$ outside $Y$ needs a distinct image adjacent to every
vertex of $T_w=N_J(w)\subseteq Y$. For every set $U$ of vertices
to be matched, Hall's criterion asks for at least $|U|$ candidate
images whose neighborhoods contain at least one $T_w$ with $w\in U$.
The family of requirements indexed by $U$ can be large. The next
lemma reduces it to a bounded subfamily after an independent
random extension.

\subsection{Compressing requirements}

For a family $\calQ$ of subsets of a finite set $Y$, say that
$S\subseteq Y$ \emph{satisfies} $\calQ$ if it contains some
member of $\calQ$. Given $r\in(0,1)$, define
\begin{equation}\label{eq:smoothed-test}
 f_{\calQ}(S)=\Prob(S\cup Z\text{ satisfies }\calQ),
 \qquad Z\sim\mu_r\text{ on }Y.
\end{equation}
A small subfamily need not be satisfied by every set that
satisfies the original family. The lemma replaces this implication
by a uniform probability bound after the extension.

\begin{lemma}[Compression after a random extension]
\label{lem:compression}
Let $Y$ be finite, $r\in(0,1)$, and $s\geq1$ an integer.
For every nonempty family
\[
 \calR\subseteq\{T\subseteq Y:1\leq|T|\leq s\}
\]
there is a nonempty subfamily $\calQ\subseteq\calR$ with
$|\calQ|\leq\lceil2r^{-s}\rceil$ such that
\begin{equation}\label{eq:compression}
 S\text{ satisfies }\calR\quad\Longrightarrow\quad
 f_{\calQ}(S)\geq1/2.
\end{equation}
\end{lemma}
\begin{proof}
Start with $\calQ=\varnothing$. While some $T\in\calR$
satisfies
\[
 \Prob(Z\text{ satisfies }\calQ\mid Z\supseteq T)<1/2,
\]
add $T$ to $\calQ$. The probability that $Z$ satisfies the
current family increases by
\[
 \Prob(Z\supseteq T,\ Z\text{ does not satisfy }\calQ)
 >\tfrac12r^{|T|}\geq\tfrac12r^s
\]
at each addition. A previously added set cannot be chosen again,
and the process makes fewer than $2r^{-s}$ additions. It makes
at least one, because $\calR$ is nonempty.

At termination the conditional probability is at least $1/2$
for every $T\in\calR$. The conditional law of $Z$ given
$Z\supseteq T$ is the law of $T\cup Z$ with a fresh Bernoulli
sample. Hence $f_{\calQ}(T)\geq1/2$. The function
$f_{\calQ}$ is increasing in $S$, proving
\eqref{eq:compression} whenever $S$ contains such a $T$.
\end{proof}

\subsection{The transfer lemma}

All logarithms below are natural. The parameter range used in both
this section and the next is
\begin{equation}\label{eq:parameters}
 0<r<1,\qquad s\in\mathbb N,\qquad
 s\leq10\log n,\qquad r^s\geq n^{-1/100}.
\end{equation}
The bounds will be uniform over this range.

\begin{proposition}[Neighborhood transfer]\label{prop:transfer}
There is an absolute integer $N_1$ such that the following holds
for every $n\geq N_1$ and every $r,s$ satisfying
\eqref{eq:parameters}. There exists
$\calD=\calD(n,r,s)\subseteq2^{E(K_n)}$ with
$\mu_r(\calD)\geq1-1/K_0$ such that, whenever
\begin{itemize}
\item $Y\subseteq[n]$ has $|Y|\leq n^{2/3}$;
\item $J$ is a labelled bipartite graph with all edges between
      $Y$ and $[n]\setminus Y$, and $E(J)\subseteq E(G)$
      for some $G\in\calD$;
\item at most $n/8$ vertices outside $Y$ have positive degree in
      $J$, and all of those degrees are at most $s$;
\end{itemize}
one has
\begin{equation}\label{eq:transfer}
 p_c(J;n)\leq r_*:=1-(1-r)^{32}\leq32r.
\end{equation}
\end{proposition}
\begin{proof}
We impose upper bounds on source sums of $f_{\calQ}$ and lower
bounds on target counts satisfying $\calQ$. Compression makes
each vertex in a set $U$ to be matched contribute at least $1/2$
to an appropriate source sum, while each target neighborhood
counted by the corresponding $\calQ$ is usable for at least one
vertex of $U$.

\paragraph{The test families.}
Use every $Y\subseteq[n]$ of size at most $n^{2/3}$ and every
nonempty family
\[
 \calQ\subseteq\{T\subseteq Y:1\leq|T|\leq s\},
 \qquad |\calQ|\leq\lceil2r^{-s}\rceil.
\]
Write $m=n-|Y|$, put $\tau=1-(1-r)^2$, and let
$x(\calQ)$ and $x_*(\calQ)$ be the probabilities that
Bernoulli-$\tau$ and Bernoulli-$r_*$ subsets of $Y$ satisfy
$\calQ$, respectively. Both are at least $r^s$.
Define $\calD$ by requiring, for every test $(Y,\calQ)$,
\begin{equation}\label{eq:source-test}
 \frac1m\sum_{w\in[n]\setminus Y}
 f_{\calQ}(N_G(w)\cap Y)\leq2x(\calQ).
\end{equation}
Also define a target event by the inequalities
\begin{equation}\label{eq:target-test}
 \frac1m\bigl|\{w\in[n]\setminus Y:
       N_{G_*}(w)\cap Y\text{ satisfies }\calQ\}\bigr|
 \geq\frac34x_*(\calQ)
\end{equation}
for all of the same tests.

\paragraph{Simultaneous concentration.}
For fixed $Y,\calQ$ and $G\sim G(n,r)$, the summands in
\eqref{eq:source-test} are independent and lie in $[0,1]$.
Their mean is $x(\calQ)$: the neighborhood sample and the
independent sample in \eqref{eq:smoothed-test} have union
density $\tau$. Their independence follows because the edges
from distinct vertices outside $Y$ into $Y$ are disjoint sets
of random variables.

The logarithm of the number of tests is at most
\begin{equation}\label{eq:test-count}
 C\bigl(n^{2/3}\log n+n^{1/100}(\log n)^2\bigr)
\end{equation}
with an absolute $C$. Indeed, choosing $Y$ costs at most
$(\lfloor n^{2/3}\rfloor+1)n^{n^{2/3}}$ possibilities.
One can then list at most $\lceil2n^{1/100}\rceil$ subsets,
each of size at most $s\leq10\log n$. This overcounts families
and gives \eqref{eq:test-count}.

For large $n$, $m\geq n/2$. Since $x(\calQ)\geq n^{-1/100}$,
Hoeffding's inequality \cite{Hoeffding63} bounds the failure
probability of \eqref{eq:source-test} for a fixed test by
$\exp(-c n^{98/100})$. The indicators in
\eqref{eq:target-test}, for $G_*\sim G(n,r_*)$, are likewise
independent, with mean $x_*(\calQ)\geq n^{-1/100}$.
The same bound, with another absolute $c>0$, applies to their
relative deviation by $1/4$.
Thus each simultaneous event has probability at least
$1-\varepsilon_n$, where uniformly in $r,s$,
\begin{equation}\label{eq:uniform-error}
 \varepsilon_n\leq
 2\exp\!\left(
 C\bigl(n^{2/3}\log n+n^{1/100}(\log n)^2\bigr)
 -c n^{98/100}\right)\longrightarrow0.
\end{equation}
Increase $N_1$ so that this error is at most
$\min\{1/K_0,1/2\}$ for $n\geq N_1$.

The relation between the two test means is important. A
Bernoulli-$r_*$ set is the union of $16$ independent
Bernoulli-$\tau$ sets. Satisfaction is increasing, so
\begin{equation}\label{eq:amplification}
 x_*(\calQ)\geq1-(1-x(\calQ))^{16}.
\end{equation}

\paragraph{Hall's criterion.}
Fix $J,Y,G$ as in the statement and a target graph $G_*$
satisfying \eqref{eq:target-test}. Keep $Y$ fixed and match
the positive-degree vertices of $J$ outside $Y$ to distinct
vertices of $[n]\setminus Y$. A candidate image $a$ is usable
for a vertex $w$ when
$N_J(w)\subseteq N_{G_*}(a)\cap Y$.

For any nonempty set $U$ of vertices to be matched, apply
Lemma~\ref{lem:compression} to
\[
 \calR=\{N_J(w):w\in U\}.
\]
It gives an admissible test $\calQ\subseteq\calR$.
Every $N_G(w)\cap Y$, for $w\in U$, satisfies $\calR$,
because $G$ contains the labelled edges of $J$. Writing
$x=x(\calQ)$, \eqref{eq:source-test} therefore yields
\begin{equation}\label{eq:hall-source}
 |U|/2\leq\sum_{w\in[n]\setminus Y}
       f_{\calQ}(N_G(w)\cap Y)\leq2xm.
\end{equation}
Repeated neighborhoods disappear from $\calR$, but their
distinct vertices still contribute separately to this sum.

Every target vertex counted in \eqref{eq:target-test} is usable
for some vertex of $U$, because $\calQ\subseteq\calR$.
If $x\leq1/16$, then
\[
 x_*(\calQ)\geq1-e^{-16x}\geq8x,
\]
so there are at least $6xm\geq|U|$ usable vertices.
If $x>1/16$, there are at least
\[
 \frac34(1-e^{-1})m\geq\frac38(1-e^{-1})n>n/8\geq|U|.
\]
Hall's theorem \cite{Hall35} gives the matching in both cases.
All choices of $\calQ$ are legitimate even though they depend
on $J,U,G$: the inequalities were established simultaneously
for every test.

The matching embeds every edge of $J$ while fixing $Y$.
Any remaining isolated vertices can be added afterward. The
target event has probability at least $1/2$, giving
$p_c(J;n)\leq r_*$. Finally, the union bound for $32$ independent
Bernoulli-$r$ samples gives $r_*\leq32r$.
\end{proof}

\section{Direct embedding lemmas}
\label{sec:direct}

Neighborhood transfer will handle edges between a small
high-degree set and its complement. Two direct estimates cover the other parts:
one for graphs with bounded maximum degree and one for small
star forests. Neither requires membership in the typical family.

\subsection{A maximum-degree bound}

\begin{proposition}\label{prop:degree-embedding}
There is an absolute integer $N_2$ such that, for $n\geq N_2$
and $r,s$ satisfying \eqref{eq:parameters}, every graph $J$
with
\[
 |V_+(J)|\leq n/4,\qquad \degen(J)\leq s,
 \qquad\Delta(J)\leq n^{2/3}
\]
and at most $n$ vertices satisfies $p_c(J;n)\leq r$.
\end{proposition}
\begin{proof}
Discard isolated vertices and assume that an edge remains.
We partition the vertices into layers that are independent sets, constructed
in reverse order. In any residual subgraph with $m'$ vertices,
the average degree is at most $2s$. At least $m'/2$ vertices
therefore have degree at most $4s$. The graph they induce has
an independent set of size at least $m'/(2(4s+1))$, by greedily
choosing a vertex and deleting it and its neighbors. Take such
a set as the last remaining layer and repeat on the residual
graph. Geometric decrease gives
\begin{equation}\label{eq:layers}
 t\leq C s\log n\leq C'(\log n)^2
\end{equation}
layers. Every vertex has at most $4s$ neighbors in earlier layers.

For each layer choose a disjoint bin in $[n]$ of size equal to
the layer size plus
\[
 b=\left\lfloor\frac{n}{4t}\right\rfloor.
\]
The bins fit, since their total size is at most $n/4+tb\leq n/2$.
Embed the layers in order. At a given step expose the edges
between the current bin and the earlier images, and seek a
matching of the layer into its bin. These edges are fresh
independent Bernoulli-$r$ variables conditional on the entire
past: all previous exposures were between earlier bins.

Write $\Delta=\Delta(J)$. Fix a subset $U$ of the current
layer with $|U|=u\geq1$. Each vertex of $U$ requires adjacency
to at most $4s$ earlier images. Each earlier image occurs in at
most $\Delta$ requirements, so the conflict graph joining
intersecting requirements has maximum degree at most $4s\Delta$.
It has an independent set of size at least
$u/(1+4s\Delta)$, whose requirements are pairwise disjoint.

For a set $W$ of $u-1$ vertices in the current bin, at least $b$
vertices of the bin remain outside $W$. The probability that
none is usable for any vertex of $U$ is at most
\begin{equation}\label{eq:hall-failure}
 \exp\!\left(-\frac{bu r^{4s}}{1+4s\Delta}\right).
\end{equation}
Indeed, restrict to the disjoint requirements just selected.
Each succeeds at a given candidate with probability at least
$r^{4s}$, independently over those requirements and over
candidates. If an empty requirement is present, the failure
event is impossible, so the bound remains valid.

Uniformly under \eqref{eq:parameters} and \eqref{eq:layers},
for sufficiently large $n$,
\begin{equation}\label{eq:hall-exponent}
 \frac{b r^{4s}}{1+4s\Delta}
 \geq c\frac{n^{1-4/100-2/3}}{(\log n)^3}
 =c\frac{n^{22/75}}{(\log n)^3}\geq5\log n.
\end{equation}
Here $b\geq c'n/(\log n)^2$ and $r^{4s}\geq n^{-4/100}$.
There are at most $n^{2u}$ choices of $U,W$ for each $u$.
Failure of Hall's condition implies such a pair with all usable
images of $U$ in $W$. By \eqref{eq:hall-failure}--\eqref{eq:hall-exponent},
the conditional failure probability at a layer is at most
\[
 \sum_{u\geq1}n^{2u}e^{-5u\log n}
 =\sum_{u\geq1}n^{-3u}=O(n^{-3}).
\]
Summing this bound over the $O((\log n)^2)$ layers shows that
the embedding succeeds with probability tending to one,
uniformly over the allowed parameters. Taking an absolute
$N_2$ so that it succeeds with probability at least $1/2$
proves the proposition.
\end{proof}

\subsection{Small star forests}

A star forest is a graph whose nontrivial components are stars.
The next argument adapts the fresh-center embedding in
\cite[Claim~6.6]{AHPT26}. Small support is imposed directly,
so it is unnecessary to divide the stars by their degrees in
advance.

\begin{proposition}\label{prop:stars}
Put $A=400e^4$. There is an absolute integer $N_3$ such that,
for every $n\geq N_3$ and every star forest $J$ on at most $n$
vertices with $|V_+(J)|\leq n/100$,
\begin{equation}\label{eq:star-threshold}
 p_c(J;n)\leq2Aq(J;n).
\end{equation}
If $q(J;n)<1/(2A)$, the stronger bound
$p_c(J;n)\leq Aq(J;n)$ holds.
\end{proposition}
\begin{proof}
The edgeless case is immediate. Put $q_J=q(J;n)>0$.
If $q_J\geq1/(2A)$, then $2Aq_J\geq1$, so
\eqref{eq:star-threshold} is automatic. We may therefore
assume $q_J<1/(2A)$ and set $p=Aq_J\in(0,1/2)$.

For $\ell\geq1$, let $m_\ell$ be the number of components
of $J$ with exactly $\ell$ leaves, ignoring isolated vertices.
If $m=m_\ell>0$, the copies on $[n]$ of $m$ disjoint
$\ell$-edge stars cover $\calF_J$. Their number is at most
\[
 C_\ell=\frac{n^{m(\ell+1)}}{(\ell!)^m m!}.
\]
This follows by counting vertex injections and dividing by
permutations of the leaves and of the stars. For $\ell=1$
there are additional automorphisms, which only decrease the
actual count. Since such a star forest exists on at most $n$
vertices, $C_\ell\geq1$. Using its cover at
$(1/(2C_\ell))^{1/(m\ell)}\in(0,1)$ yields
\begin{equation}\label{eq:star-count}
 q_J^\ell\geq\frac{\ell!}{n^{\ell+1}}(m!/2)^{1/m}
 \geq\frac{\ell!}{n^{\ell+1}}\frac{m_\ell}{2e}.
\end{equation}
The last step uses $m!\geq(m/e)^m$ and $2^{1/m}\leq2$.

Partition $[n]$ into a center bin of size $\lfloor n/2\rfloor$
and a leaf bin containing the remaining vertices. Process the
components of $J$ one at a time. For a star with $\ell$ leaves,
test unused centers successively until one has at least $\ell$
neighbors among the unused leaf vertices. Use that center and
any $\ell$ such neighbors, and never test that center again.
Failed centers are also discarded.

Throughout the process, at least $b'=\lfloor n/3\rfloor$
leaf vertices remain. The unused leaf set is determined by
previously tested center rows, and a fresh center row has not
been exposed. Thus each trial, conditional on the past, has
success probability at least
\begin{equation}\label{eq:star-success}
 u_\ell=\Prob(\Bin(b',p)\geq\ell).
\end{equation}
For the purpose of counting trials, append infinitely many
auxiliary centers with fresh independent rows to the same leaf
bin. This enlarged process agrees with the actual process until
the original center bin is exhausted. The number $T$ of trials
needed to finish all stars in the enlarged process satisfies
\begin{equation}\label{eq:expected-trials}
 \E T\leq\sum_{\ell:m_\ell>0}\frac{m_\ell}{u_\ell},
\end{equation}
because the conditional expected waiting time for each success
is at most $1/u_\ell$.

If $\mu=b'p\geq2\ell$, the second-moment bound
$\E[\Bin(b',p)^2]\leq\mu^2+\mu$ and the Paley--Zygmund
inequality give
\[
 u_\ell\geq\frac14\frac{\mu^2}{\mu^2+\mu}\geq\frac16,
\]
since $\mu\geq2$. Otherwise $b'p<2\ell$. For every occupied
star size, $\ell\leq n/100$, and for large $n$,
$b'-\ell+1\geq n/4$. As $p<1/2$, the inequality
$\log(1-p)\geq-2p$ implies
\begin{align}
 u_\ell
 &\geq\binom{b'}{\ell}p^\ell(1-p)^{b'-\ell}
 \geq\frac{(n/4)^\ell p^\ell}{\ell!}e^{-4\ell}\notag\\
 &\geq\left(\frac{A}{4e^4}\right)^\ell
        \frac{m_\ell}{2e n}
 =100^\ell\frac{m_\ell}{2e n},
 \label{eq:star-small-mean}
\end{align}
where the last inequality uses \eqref{eq:star-count}.

There are at most $n/200$ nontrivial components. Combining
the two bounds in \eqref{eq:expected-trials} gives
\[
 \E T\leq6\frac{n}{200}+2e n\sum_{\ell\geq1}100^{-\ell}
 =\left(\frac3{100}+\frac{2e}{99}\right)n<\frac n8.
\]
For large $n$, Markov's inequality shows that
$T\leq\lfloor n/2\rfloor$ with probability at least $1/2$.
On this event all centers used belong to the actual random graph,
so the construction embeds $J$ in $G(n,p)$. This proves
$p_c(J;n)\leq Aq_J$ in the remaining case.
\end{proof}

\section{Constructing the fixed partition}
\label{sec:assembly}

We now prove Theorem~\ref{thm:main}. The high-degeneracy case
first uses the typical-family cover and isolates a small set $Y$
so that every vertex outside $Y$ has few neighbors in $Y$ and
bounded total degree. Neighborhood transfer handles edges across
this cut, while the maximum-degree estimate handles edges inside
and outside $Y$ after further splitting. The bounded-degeneracy
case uses only small star forests. The partition choices in both
cases are deterministic.

\begin{proof}[Proof of Theorem~\ref{thm:main}]
Fix
\begin{equation}\label{eq:constants}
 M=B=1000,\qquad k=2K_0M(B+1)^2.
\end{equation}
Choose an absolute integer $n_0\geq2$ large enough for
Propositions~\ref{prop:transfer}, \ref{prop:degree-embedding},
and~\ref{prop:stars}, and so that, for every $n\geq n_0$,
\begin{equation}\label{eq:deterministic-cutoffs}
 \frac{8\log n}{\log2}\sqrt n\leq n^{2/3},\qquad
 n^{2/3}\leq n/4,\qquad
 2\lceil n/B\rceil\leq n/100.
\end{equation}
All these conditions have absolute cutoffs. Set
\[
 L=\max\{2A,64,n_0^2\}.
\]
If $H$ is edgeless, use only empty pieces. Otherwise write
$q=q(H;n)>0$ and $d=\degen(H)\geq1$.
When $n<n_0$, Lemma~\ref{lem:q-lower} gives
$Lq\geq n_0^2/(n(n-1))\geq1$. When $q\geq1/(2A)$,
we also have $Lq\geq1$. In either case put all edges in one
piece and use empty pieces for the rest. Henceforth
\begin{equation}\label{eq:nontrivial-range}
 n\geq n_0,\qquad q<1/(2A).
\end{equation}

\paragraph{Large degeneracy: $d\geq M$.}
Set
\[
 r=2q,\qquad s=\lceil2d/M\rceil.
\]
By Lemma~\ref{lem:q-lower}, $r\geq n^{-2/d}$, whereas
\eqref{eq:nontrivial-range} gives $r<1/A<1/2$. Consequently
\begin{equation}\label{eq:d-upper}
 d<\frac{2\log n}{\log2},\qquad s\leq3d/M,
 \qquad r^s\geq n^{-6/M}\geq n^{-1/100}.
\end{equation}
These bounds also give $s\leq10\log n$, so
\eqref{eq:parameters} holds.

Take the family $\calD$ from Proposition~\ref{prop:transfer}.
Corollary~\ref{cor:typical-cover} supplies a copy of $H$ covered
by $K_0$ members $G_1,\ldots,G_{K_0}$ of $\calD$.
Fix those graphs and the labelled copy, and use this copy to
label $H$ on a subset of $[n]$. Assign each edge of $H$ to
one $G_i$ containing it. This gives $K_0$ initial edge pieces,
each contained in its specified typical graph.

We separate the high-degree vertices, following the obstacle
identified in \cite[Section~8]{AHPT26}. Initially let
\[
 Y_0=\{v\in V(H):\deg_H(v)>\sqrt n\}.
\]
By \eqref{eq:edge-bound}, $|Y_0|\leq2d\sqrt n$.
Starting from $Y_0$, add an outside vertex whenever it has more
than $2d$ neighbors in the current set, and stop when no such
vertex remains. Let $Y$ be the resulting set and let
$a=|Y\setminus Y_0|$. Each added vertex contributes more than
$2d$ edges to earlier vertices of the set; no edge is counted
twice. Thus
\[
 2da\leq|E(H[Y])|\leq d(|Y_0|+a),
\]
giving $|Y|\leq2|Y_0|\leq4d\sqrt n\leq n^{2/3}$ by
\eqref{eq:d-upper} and \eqref{eq:deterministic-cutoffs}.
The conclusion also holds if $Y_0$ is empty, in which case no
vertex is added. At termination every $w\notin Y$ satisfies
\begin{equation}\label{eq:outside-bounds}
 \deg_H(w)\leq\sqrt n,\qquad |N_H(w)\cap Y|\leq2d.
\end{equation}

Partition $V(H)\setminus Y$ into $B$ possibly empty blocks
of size at most $\lceil n/B\rceil$. Subdivide each initial
edge piece as follows.

\emph{Edges inside $Y$.}
Apply Lemma~\ref{lem:split} with $m=M$ to obtain $M$ pieces
of degeneracy at most $\lceil d/M\rceil\leq s$.
Each is supported on at most $|Y|\leq n^{2/3}\leq n/4$
vertices and has maximum degree at most $n^{2/3}$.
Proposition~\ref{prop:degree-embedding} gives threshold at most $r$.

\emph{Edges outside $Y$.}
First partition by the unordered pair of endpoint blocks,
allowing both endpoints in the same block. There are
$B(B+1)/2$ such pairs. Apply Lemma~\ref{lem:split} with
$m=M$ to each resulting graph. The pieces have degeneracy
at most $s$, maximum degree at most $\sqrt n$, and support
size at most $2\lceil n/B\rceil\leq n/100$.
Proposition~\ref{prop:degree-embedding} again gives threshold
at most $r$.

\emph{Edges between $Y$ and its complement.}
First partition by the block of the endpoint outside $Y$.
At each such endpoint, distribute its edges in the initial
piece among $M$ groups of size at most $\lceil2d/M\rceil=s$,
using \eqref{eq:outside-bounds}. This assigns every cut edge
once. Each resulting bipartite piece is still contained in its
specified $G_i\in\calD$, has at most $\lceil n/B\rceil\leq n/8$
positive-degree vertices outside $Y$, and has all those degrees
at most $s$. Proposition~\ref{prop:transfer} gives threshold
at most $r_*\leq32r=64q$.

All pieces therefore have threshold at most $64q\leq Lq$.
Their number is at most
\[
 K_0M\left(1+\frac{B(B+1)}2+B\right)
 =\frac{K_0M(B+1)(B+2)}2\leq k.
\]

\paragraph{Bounded degeneracy: $1\leq d<M$.}
Lemma~\ref{lem:split}, with $m=d$, partitions $E(H)$ into
$d$ graphs of degeneracy at most one, hence into forests.
Indeed, the last vertex of a cycle in the certifying order
would have two preceding neighbors.

Root every nontrivial tree. Color each parent--child edge by
the parity of the depth of its parent. Each color class is a
star forest: its centers occur at one depth parity, and their
children have the other parity and no other incident edge of
that color. This is the parity splitting used in
\cite[Claim~6.2]{AHPT26}.

Partition $V(H)$ into $B$ blocks of size at most
$\lceil n/B\rceil$, and split each star forest by unordered
endpoint-block pairs. Every resulting piece $J$ remains a star
forest and has
$|V_+(J)|\leq2\lceil n/B\rceil\leq n/100$.
By Proposition~\ref{prop:stars} and \eqref{eq:monotonicity},
\[
 p_c(J;n)\leq2Aq(J;n)\leq2Aq(H;n)\leq Lq.
\]
The number of pieces is at most
$2d\,B(B+1)/2=dB(B+1)\leq k$.

Finally, add empty pieces if necessary. In the large-degeneracy
case, pull the labelled edge pieces back to the original $H$.
The typical graphs, the copy of $H$, all vertex orders, and
all assignments have been fixed using only $H,n$; none depends
on the random host used to test a piece's threshold. This
completes the proof.
\end{proof}

\bibliographystyle{amsplain}
\bibliography{references}
\end{document}
